\section{Scope and Sharpness of the Hypotheses}
\label{sec:scope-and-sharpness-of-the-hypotheses}

The hypotheses in \eqref{eq:main-assumptions} include centered,
variance-normalized Laplace densities, despite their cusp at the origin.
They also include variance-normalized densities proportional to $e^{-x^4}$.
Symmetry is not needed in any step of the proof.

For examples allowing both asymmetry and heavy tails, take a density
proportional to
\begin{equation}
 (1+y^2)^{-(\tau+1)/2}
                           e^{\beta\arctan y},
 \qquad \tau>2,\quad \beta\in\R,
 \label{eq:asymmetric-heavy-tail-example}
\end{equation}
and center and rescale to variance one.
Its location score before centering and rescaling is
\[
 \frac{(\tau+1)y-\beta}{1+y^2}.
\]
Both this score and its product with $y$ are bounded.
The density is positive and smooth.
It has a finite second moment.
Centering and rescaling preserve finiteness of the location and scale
Fisher information.
So the standardized density satisfies \eqref{eq:main-assumptions}.
For \(2<\tau\le4\), its fourth moment is infinite.
When \(\beta\ne0\), the two tails have unequal leading constants.
So centering does not make the density even.
The case \(\beta=0\) is a rescaled Student $t$ family.
These examples show that neither evenness nor a fourth moment is implicit in
the theorem.

The Gamma example used to illustrate the rotational Fourier defect in
\cref{sec:variance-and-affine-normalization} has a centered density that vanishes on
a half-line.
Thus it does not satisfy the positivity hypothesis of
\cref{thm:measure-equivalence}.
A nonzero rotational Fourier defect establishes a local Hellinger lower bound.
Alone, this does not establish the general measure-equivalence theorem for a
density with restricted support.
\Cref{thm:restricted-support-rigidity} supplies the missing classification:
for every standardized Gamma shape, only coordinate permutations give
equivalence under homogeneous bounded invertible linear mixing.
The same theorem classifies all bounded coordinate laws, including
disconnected supports and laws without densities.
\Cref{cor:restricted-gamma-affine} treats the distinct affine Gamma
problem, where support-preserving scales do allow equivalence.

The second moment has a specific role beyond differentiating the expectation
of the product of centered complex exponentials.
It makes every \(\ell^2\) row into a convergent random series.
It supplies uniform integrability of squares of linear forms.
It gives the coefficient-subsequence argument with an independent Gaussian
summand in \cref{lem:unit-row-rigidity}.
So the first-moment ODE result is stronger in its moment scope than the
arbitrary-operator equivalence theorem.
Extending the latter to infinite-variance coordinates requires a different
domain of admissible rows and a corresponding operator criterion.
For example, for symmetric nondegenerate \(\alpha\)-stable coordinates,
\(0<\alpha<2\), convergence of a row series requires
$\sum_j|a_j|^\alpha<\infty$: its partial-sum characteristic
functions are
\[
 \exp\left(-c|t|^\alpha\sum_{j\le n}|a_j|^\alpha\right).
\]
An arbitrary \(\ell^2\) row need not meet that requirement.
Using bounded test functions does not remove this condition on admissible rows.

Under the density positivity hypothesis, the Hellinger path estimates give $L^1$ convergence and uniform
integrability of likelihood ratios.
The finite real Fourier sums are bounded in each coordinate.
Their dimension-free bounds are the $L^2(\mu_n)$ and generator
estimates proved above.
No variance bound for their aggregate under an arbitrary mixed output law is
needed in the tail-symmetry argument.
